\documentclass[12 pt, letterpaper]{hmcpset}
\usepackage{lin-alg-preamble}

\name{Jayden Thadani}
\class{MATH 171 --- Abstract Algebra I}
\assignment{Homework 8}
\duedate{16th February 2024}

\begin{document}

\begin{problem}[Dummit \& Foote \S2.1 \#6]
	Let $G$ be an abelian group. Prove that $\{ g \in G | \lvert g \rvert < \infty \}$ is a subgroup of $G$ (called the \textit{torsion subgroup} of $G$). Give an example where this set is not a subgroup when $G$ is non-abelian.
\end{problem}

\begin{solution} 
	Suppose $G$ is an arbitrary abelian group. Let $H$ denote the set of all elements of finite order. We want to show that $H$ is a subgroup of $G$. Note that we only need to show that for all $x, y \in H$ we have $xy^{-1} \in G$ (and that $H$ is nonempty). 

	We have trivially that $H$ is nonempty since $1_G$ has finite order. Suppose $x, y \in H$. Specifically, suppose $\lvert x \rvert = m$ and $\lvert y \rvert = n$ with $m, n \in \mathbb{N}$. Then note that
	\[
		\begin{split}
			(xy^{-1})^{mn} & = \underbrace{(xy^{-1}) \dots (xy^{-1})}_{mn \text{ times}} \\
				       & = x^{mn} (y^{-1})^{mn} \\
				       & = (x^{m})^n (y^n)^{-m} \\
				       & = 1_G.
		\end{split}
	\]
	Thus, $\lvert xy^{-1} \rvert \le mn$ and $xy^{-1} \in H$ as desired. That is, $H$ is a subgroup --- the torsion subgroup of $G$.

	Consider $G = \mathrm{GL}_2(\mathbb{R})$. Then $H$ contains $\begin{bmatrix}
			0 & 1 \\
			1 & 0
		\end{bmatrix}$ and $\begin{bmatrix}
			0 & 2 \\
			1/2 & 0
		\end{bmatrix}$ which are both of order $2$. However, $H$ does not contain their product
		\[
			\begin{bmatrix}
				0 & 1 \\
				1 & 0
			\end{bmatrix}
			\begin{bmatrix}
				0 & 2 \\
				1/2 & 0 \\
			\end{bmatrix} = 
			\begin{bmatrix}
				1/2 & 0 \\
				0 & 2
			\end{bmatrix}
		\]
		since
		\[
			\begin{bmatrix}
				1/2 & 0 \\
				0 & 2
			\end{bmatrix}^n =
			\begin{bmatrix}
				1/2^n & 0 \\
				0 & 2^n
			\end{bmatrix}
		\]
		which is never $I$ for any $n \in \mathbb{N}$, and thus $\begin{bmatrix}
			1/2 & 0 \\
			0 & 2
		\end{bmatrix}$ is not of finite order.
\end{solution}

\pagebreak

\begin{problem}[Dummit \& Foote \S2.1 \#8]
	Let $H$ and $K$ be subgroups of $G$. Prove that $H \cup K$ is a subgroup if and only if either $H \subset K$ or $K \subset H$.
\end{problem}

\begin{solution}
	First we will show that if $H \subset K$ or $K \subset H$  then $H \cup K$ is a subgroup of $G$. If either of the assumptions is true, then we have $H \cup K = K$ or $H \cup K = H$ respectively. Since they are already subgroups, we have that $H \cup K$ is a subgroup. 

	Now, suppose $H \not \subset K$ and $K \not \subset H$. Then we must have some $k \in K \setminus H$ and some $h \in H \setminus K$. We claim that though we have $h, k \in H \cup K$, we do not have $hk \in H \cup K$. We will show this by contradiction. 

	Suppose we had $hk \in H \cup K$ and further, without loss of generality, suppose that we had $hk \in H$. Then, since $h^{-1} \in H$ we would have to have $h^{-1}hk = k \in H$, contradicting that $k \in K \setminus H$. Thus our assumption that $hk \in H \cup K$ must be false.
\end{solution}

\pagebreak

\begin{problem}[Dummit \& Foote \S2.1 \#10]
	\begin{enumerate}
		\item Prove that if $H$ and $K$ are subgroups of $G$ then so is their intersection $H \cap K$.
		\item Prove that the intersection of an arbitrary nonempty collection of subgroups of $G$ is again a subgroup of $G$ (do not assume the collection is countable).
	\end{enumerate}
\end{problem}

\begin{solution}
	\begin{enumerate}
		\item Note that $H \cap K$ is non-empty since $H$ and $K$ both contain the identity $1_G$. Suppose $g_1, g_2 \in H \cap K$. Then $g_1, g_2$ are in both, $H$ and $K$. Thus, by the subgroup criterion, $g_1 g_2^{-1}$ is in both $H$ and $K$. Thus, $g_1 g_2^{-1} \in H \cap K$, giving us that it is a subgroup.
		\item Suppose we have $\{ H_i \}$, a collection of subgroups of $G$, indexed by $i \in I$. Let $H = \bigcap_{i \in I} H_i$. Notice that $H$ is nonempty since $1_G \in H_i$ for all $i \in I$. Further, suppose that we have $h_1, h_2 \in H$. Then we have $h_1, h_2 \in H_i$ for all $i \in I$. By the subgroup criterion we have $h_1 h_2^{-1} \in H_i$ for each $i \in I$ and thus, we have $h_1 h_2^{-1} \in H$ which by the subgroup criterion gives us that $H$ is a subgroup.
	\end{enumerate}
\end{solution}

\pagebreak

\begin{problem}[Dummit \& Foote \S2.1 \#15]
	Let $H_1 \le H_2 \le \dots$ be an ascending chain of subgroups of $G$. Prove that $\bigcup^\infty_{i = 1} H_i$ is a subgroup of $G$.
\end{problem}

\begin{solution}
	Let $H = \bigcup_{i = 1}^\infty$. Notice that $H$ is clearly nonempty since it contains all of these nonempty $H_i$. Suppose $h_1, h_2 \in H$ with $h_1 \in H_j$ and $h_2 \in H_k$. Without loss of generality, assume $k \ge j$. Then $H_j \subset H_k$ and thus, $h_1 \in H_k$. Thus, by the subgroup criterion, we have $h_1 h_2^{-1} \in H_k$, and thus, $h_1 h_2^{-1} \in H$, giving us that $H$ is a subgroup.
\end{solution}

\end{document}
